% 本文件由 LaTeX 排版 Agent 根据有效 translation.json 自动生成。
% author=John Chrysostom; work_id=1913; title=论反对泄露弟兄们的错误

\chapter{论反对泄露弟兄们的错误}

% structure: p0001 source=h1 target=omitted reason=page-title-already-rendered-by-parent

1. \Emph{我认为}你们是有福的，亲爱的，因你们怀着热忱涌向圣父的家。因为从这热忱中，我有理由对你们灵魂的健康也抱有自信；的确，教会的学校是一间绝妙的手术室——不是为身体，而是为灵魂的手术室。因为那是属神的，矫正的不是肉体的创伤，而是心灵的错误；对这些错误和创伤，药物就是圣言。这药物不是由地上生长的草药调和，而是由自天而降的言语配制——不是医师的手，而是宣讲者的唇舌分施的。因此它持久不变；既不被时间消磨其效力，也不被任何疾病的强度击败。因为医师的药物确实有这两种缺陷：它们新鲜时显示其固有力量，但时间一长，就像衰老的身体一样变得虚弱；而且，恶疾的顽固也常常使它们束手无策，因为它们是人的产物。而神圣的药物却并非如此；经过很长时间之后，它仍然保有它所有内在的效力。至少从\Person{梅瑟}时代起（因为圣经的开端可追溯至此），它已治愈了如此众多的人；它不仅没有失去其特有的力量，也从未有任何疾病能战胜它。这药物无法用银钱购得；但谁若展现了真诚的意向与心志，便可将其全部带走。为此，富人和穷人都同样获得这治疗过程的益处。因为凡需支付银钱的地方，固然是财力雄厚者获益；但穷人常常因无法负担配制药物的费用而空手离去。然而在此，既不可能支付银钱，却需要显示信德和善意，那怀着乐意的心付出这些的人，正是获益最多的人；因为这些正是治疗所需的代价。富人和穷人同样获益；或者更确切地说，获益并不相同，穷人往往带着更多的收获离去。这究竟是为什么？因为富人先被许多挂虑占据，有财富带来的骄傲和虚荣，以粗心和懒散为伴，聆听圣经的药物时并不十分专注，也不十分恳切；而穷人远离了奢侈、贪食和怠惰，将全部时间用在手工和正直的劳动上，由此为灵魂收集了许多智慧，从而变得更加专注、不松懈，习惯于以更仔细的留心思考所说的一切；因此，既然他所付的代价更高，他离开时所得的益处也更大。\par

2. 我说这一切，并非绝对指责那些富有的人，也不是无条件地赞美穷人：因为财富并非邪恶，但滥用财富是邪恶；贫穷并非德行，但善用贫穷是德行。那位在\Person{拉匝禄}时代的富人受惩罚并非因为他富有，而是因为他残酷无情。那位安息在\Person{亚巴郎}怀里的穷人受称赞，并非因为他贫穷，而是因为他以感恩的心承受了他的贫穷。\par

因为事物——（现在请仔细听这句话；因为它足以在你们心中植入足够的宗教知识，并驱逐一切不健全的推理，使你们对事物的真理有正确的判断）——事物中，有些本性上是道德善的，有些则相反；有些既不善也不恶，处于中间位置。虔敬在本性上是善事，不虔敬是恶事；德行是善事，恶行是恶事；但财富和贫穷本身既不是前者也不是后者；它们因使用者的意志而成为前者或后者。因为如果你用你的财富行了仁爱之事，那事物就成为你善的基础；但如果你用来掠夺、强取和傲慢，你就将它用到了完全相反的方面；但这不应归咎于财富，而应归咎于那将财富用于傲慢的人。同样，我们也可以说贫穷：如果你以感恩的心向主高尚地承受它，你所做的就成为你获得荣冠的原因和基础；但如果你为此亵渎你的造物主，并因祂的眷顾而指责祂，那么你又将这事物用于邪恶的目的。但正如在那种情况下，财富并非贪婪的根源，而是那滥用财富的人，同样，在这里我们也不应将亵渎的罪责归于贫穷，而应归于那不愿以清醒精神承受此事的人。因为在每件事上，赞美和责备都取决于我们自己的意志和选择。财富是善，但并非绝对，而只是对那没有犯罪的人而言；同样，贫穷是恶，但并非绝对，而只是在恶人的口中，因为他不知足，因为他亵渎，因为他愤慨，因为他指责那创造他的主。\par

3. 因此，我们不要指责财富，也不要绝对地辱骂贫穷，而要指责那些不愿善用它们的人；因为事物本身是中性的。但正如我所说（因为回到前一个主题是好的），富人和穷人都以同样的勇气和自由享受这里所施予的药物的益处；而穷人往往更恳切。因为这些药物的卓越之处不仅在于它们治愈灵魂，它们的效力不因时间而消减，它们不被任何疾病击败，益处免费公开提供，治疗对富人和穷人平等——它们还有另一个不亚于这些优点的品质。请问是什么？那就是我们不会公开揭露那些来到这手术室的人。因为去外部世界手术室的人，有许多人检查他们的伤口，除非医生先揭开疮口，否则他不会敷药；但这里并非如此，我们虽看到无数病人，却以隐蔽的方式对他们进行治疗。因为我们不是通过将犯罪者公之于众来宣扬他们所犯的罪；而是提出我们的教导，作为对所有人共同的教导，完全留给听者的良心；这样每个人都可以从所说的话中汲取适合自己伤口的药物。因为教导之言从讲者的舌头发出来，包含着对邪恶的控告、对美德的赞扬、对放荡的指责、对贞洁的称赞、对骄傲的谴责、对温良的赞美，如同一种混合了多种成分的药物；每个听者从中汲取适合自己且有益的部分。因此，话语公开发出，进入各人的良心，秘密地提供它所带来的治疗，并且在疾病被揭露之前，往往已恢复了健康。\par

4. 你们昨天一定听到了我是如何赞扬祈祷的力量，如何责备那些心不在焉祈祷的人；却没有公开揭露他们中的任何一个。那些自觉恳切的人接受了那对祈祷的称赞，并通过赞美变得更加恳切；而那些自觉懈怠的人，则接受了责备，并摆脱了粗心。但无论是这些人还是那些人，我们都不认识；这种无知对两者都有益——我现在告诉你们如何。那位听到对祈祷的称赞并自觉恳切的人，如果他有众多见证人见证称赞，就会滑向骄傲；但事实上，他秘密地接受了称赞，就远离了一切傲慢。另一方面，那位自觉懈怠的人，听到指责后，因指责而变得更好，因为他没有一个人见证指责；这对他有非同寻常的益处。因为由于担心俗众的看法，只要我们以为自己能在恶行中隐藏，我们就会努力变得更好；但当我们对所有人变得臭名昭著，失去了隐藏所带来的安慰时，我们就变得更加无耻和松懈。正如伤口如果不去包扎，经常暴露在冷空气中就会更加疼痛；同样，灵魂犯罪后，如果在众人面前因做错事而受责备，就会因此变得更加无耻。为了避免这种情况，话语暗中对你施以药物。为使你们明白这种暗中处理的好处有多大，请听基督所说的话。\BibleQuote{如果你的弟兄得罪了你，要劝戒他。}祂没有说\BibleQuote{在他和整个城镇之间}，也没有说\BibleQuote{在你和全体人民之间}，而是说\BibleQuote{只在你和他之间。}祂说，要指责没有见证，这样转向改过就容易接受。毫无疑问，劝告不公开是一件大好事。良心足够了，那不可腐蚀的审判官足够了。指责犯罪者的不是你，而是他自己的良心（那控告者更严厉），你对所犯之罪的了解也不如他精确。因此，不要通过揭露犯罪者而加重伤口；要在没有见证的情况下给予劝告。我们现在所做的正是这个——保禄所做的也是如此，他在控诉格林多人中犯罪者时没有引用见证。请听他是怎么说的：\BibleQuote{因此，弟兄们，我把这些比喻应用在我和阿颇罗身上。}其实，分裂人民、分裂教会的既不是他本人也不是阿颇罗；但他仍然隐藏了指责，正如用面具一样，用他自己和阿颇罗的名字掩盖被告的面孔，使他们有能力改过自新。又说：\BibleQuote{免得我来到时，天主使我谦卑，而我必须哀悼许多先前犯罪、没有悔改他们所作的不洁、淫乱和放荡的人。}请看，这里他又不具体地提到那些犯罪的人，以免公开指控使犯罪者的灵魂更加无耻。因此，正如我们以如此的体谅来施予责备，你们也要以全部认真接受改正；并留心听所说的话。\par

5．我们昨日向你们谈论了祈祷的能力。我指出了魔鬼是如何潜伏伺机，他真是个欺骗者。因为他看到从祈祷中我们获得极大的益处，因此他就最猛烈地攻击我们，为要使我们失去防卫，并使我们空手回家。正如在法官面前，那些围绕着法官的法庭侍卫，对那些前来求见的人怀有敌意时，就用棍棒将他们驱赶到远处，阻止他们靠近并流泪哀求以获取怜悯；同样，魔鬼看见我们来到审判者面前时，也借着我们自己的懈怠，而非用棍棒，将我们驱赶到远处。因为他知道，他知道得很清楚，如果人们以清醒的精神来到他面前，承认所犯的罪过，并以热切的灵魂悲哀痛悔，他们将会得到完全的宽恕而离去；因为天主爱人；因此他先发制人，阻止我们亲近，为使我们得不到任何所需的东西。但是法官的士兵是强力地吓跑那些前来的人；而他却不用强制手段，只是欺骗我们，使我们陷入安逸之中。因此我们甚至不配得到宽容，因为我们自愿剥夺了自己的美善。恳切的祈祷是理智和灵魂的光——一种不灭而永恒的光。为此，他将无数想像的垃圾堆投入我们心中；那些我们从未想象过的东西，在祈祷的那一刻聚集起来，倾注到我们的灵魂中。正如狂风常从相反方向猛烈吹来，在点燃之际就熄灭灯的火焰；同样，魔鬼看见我们祈祷的火焰被点燃时，就从不计其数的思虑之风吹动它，四面八方地吹，直到熄灭那光为止。但是，那些点燃灯的人所做的事，我们也要这样做。他们做什么呢？当他们看见强风袭来，就用手指堵住灯的开口，阻挡风的进入。因为只要他从外面攻击，我们就能抵挡他；但当我们向他敞开心灵的门户，接纳敌人进入里面后，我们就再也不能稍有抵挡了；反而完全熄灭了记忆，如同冒烟的灯一般，他让我们的口空说空洞言词。但是，正如他们用手指堵住灯的开口，我们也当把思虑放在心灵上：让我们对邪恶之神关闭入口，以免他熄灭我们祈祷的光。请记住这两个比喻，一个是关于士兵和法官的，另一个是关于灯的。我们正是为此目的向你们提出这些比喻；这些是我们所熟悉的，我们生活其中的，为使我们离开这里回家后，能从日常熟悉的事物中回想起所说的话。\par

6．祈祷是一件坚固的盔甲和极大的保障。你们昨日听到了三个青年虽被捆绑，却摧毁了火的力量；他们践踏火焰，克服了火窑，战胜了火元素的作用。今日再听一听那高贵伟大的依撒格如何借着祈祷克服了身体的本性本身。他们摧毁了火的力量，而这个人今日解开了无能本性的束缚。请学习他是如何成就的。\Quote{依撒格}，经上说，\Quote{为他的妻子祈求，因为她不生育。}今天已向你们读了这段；昨天讲道是关于祈祷，今日再次有祈祷能力的显明。请看圣神的恩宠如何安排，使今天所读的与昨天所说的和谐一致。\Quote{依撒格}，经上说，\Quote{为他的妻子黎贝加祈求，因为她不生育。}这首先值得探究，她为何不生育。她的生活是可赞美的，充满极大的贞节——她本人和她的丈夫都是如此。我们不能抓住那些义人的生活，说不生育是罪的惩罚。不仅她本人不生育，而且他的母亲撒辣，那位生了依撒格的女人，也不生育；不仅他的母亲和妻子不生育，而且他的儿媳，雅各伯的妻子辣黑耳，也是如此。这一群不孕者有什么含义呢？他们都是义人，都生活在德行中，都得了天主的见证。因为关于他们，主曾说：\Quote{我是亚巴辣罕的天主，依撒格的天主，雅各伯的天主。}关于同样的人，保禄也这样说：\Quote{因此，天主不以称自己为他们的天主为耻。}\BibleRef{希 11:16} 在新约中，有许多对他们的赞美；在旧约中，有许多对他们的颂扬。他们各方面都光明显耀，但他们的妻子都不生育，而且一直不生育直到年老。因此，当你看到夫妻生活在德行中，看到他们蒙天主喜爱，关心虔诚，却遭受不孕之症时，不要以为不孕是对罪的惩罚。因为天主的安排有很多理由，而对我们来说是不可解释的；对这一切我们都必须衷心感恩，并认为只有那些生活在邪恶中的人才是可悲的，而不是那些没有子女的人。天主常常适宜地这样做，尽管我们不知道事件的缘由。为此，在任何情况下，我们都应该赞叹他的智慧，并荣耀他不可言喻的爱人之情。\par

7. 实在，这一考量足以在道德品格上训导我们，但也有必要说明那些妇女不生育的原因。原因何在？是为了当你看到童贞玛利亚生下我们共同的主时，不至于不信。因此，要在不生育的胎中操练你的心智；以便当你看到那因天主的恩宠而残废、被束缚的胎，却开启孕育子女时，你听到童贞女生产，就不会惊奇，或者即使惊奇也仍惊叹，但不要对奇迹不信。当犹太人对你说：\Quote{童贞女如何生产？}你就对他说：\Quote{那因年老而衰弱不生育的，又如何生产？}那时有两种障碍：年龄的不适时和本性的无用；但就童贞女而言，只有一种障碍：没有婚配。因此，不生育的为童贞女预备了道路。为让你知道，正是为此，不生育的先例在先，为使童贞女的生产被相信，请听加俾额尔对她说的话——因为当他前来对她说：\BibleQuote{你将怀孕生子，要给他起名叫耶稣；}\BibleRef{路 1:31}童贞女惊奇而诧异，说：\Quote{这事如何成就在我身上，因为我不认识男人？}天使怎么说？\BibleQuote{圣神要临于你。}\BibleRef{路 1:35}他说，不要寻求本性的顺序，因为所发生的事超乎本性；不要寻求婚姻和生产的阵痛，因为生产的方式过于婚姻的尊贵。她说：\Quote{这事如何成就，因为我不认识丈夫？}实在，正因为你不认识丈夫，这事才会成就。因为若你认识丈夫，你就不配服务于这一职分。因此，你因不信的理由，反而应当相信。你也不配服务于这一职分，并非因为婚姻是邪恶，而是因为童贞更优越；并且按理，主的进入应比我们更尊贵；因它是君王的，君王通过更尊贵的途径进入。祂必须在出生上既与我们相同，又与我们相异。因此这两方面都得以安排。\par

因为从胎中出生，在我们看来是共同的；但无婚姻而出生，却是一件超越我们水平的事。而在腹中怀胎、成孕，属于人性；但怀孕不经性交，则过于人性尊贵。为此，这两件事都发生了，为让你认识那诞生者的超越性与和你的相似性。\par

8. 请思考这一切所做之事的智慧。超越性并未损害与我们的相似和亲缘，与我们的亲缘也未使超越性暗淡；反之，两者都在所有环境中显示；一方完全拥有我们的状况，另一方则拥有与我们相比的差异。但正如我所说，为此不生育的先例在先，为使童贞女的生产被相信，使她能被引导相信她从天使听到的许诺和保证。天使说：\BibleQuote{圣神要临于你，至高者的德能要庇荫你}\BibleRef{路 1:35}——这样，他说，你就能承受。不要注视大地；操作将从天上而来。所发生的事是圣神的恩宠；请不要询问本性和婚姻的规律。但由于那些话对她太高深，他愿意也提供另一种证明。但请你观察，不生育的如何引导她相信此事。因为那证明对童贞女的智力太高深，请听他是如何将他的话也降低到低层次的事物，通过可感觉的事实引导她。因为他说：\BibleQuote{看，你的亲戚依撒伯尔——她也曾在老年怀了男胎；现在已是第六个月了，她原是素称不生育的。}\BibleRef{路 1:36}你看见吗？不生育的是为了童贞女的缘故。因为他引述她亲戚的生产是为了什么目的？他说\Quote{在老年}是为了什么目的？他加上\Quote{原是素称不生育的}是为了什么目的？显然是为了通过这些一切诱导她相信喜讯。为此，他既提到年龄，也提到本性的残疾；为此，他也等待受孕后所经过的时间；因为他并非从开始就立即向她报告喜讯，而是等到不生育的已过了六个月，使怀胎的隆起从此成为怀孕的保证，并为受孕产生无可争议的证据。请再看加俾额尔的智慧。因为他既未提到撒辣，也未提到黎贝加，也未提到辣黑耳；虽然她们也曾不生育，她们也曾年老，所发生的事也是奇迹；但那些故事是古代的。而新近的、当代的、发生在我们这一代的事，比古老的事更容易诱导我们相信奇迹。为此，他放下那些妇女，让她从她的亲戚依撒伯尔本人理解将要临到她的事，于是提出了这事；为从她引导她到她自己——那最可敬畏、最尊贵的生产。因为不生育的生产介于我们和主的生产之间：确实比童贞女的低，但比我们的高。为此，他借着依撒伯尔居中，如同借着某种桥梁，将童贞女的思想从符合本性的生产，提升到超乎本性的生产。\par

9. 我原想说得更多，教给你们黎贝加和辣黑耳不生育的其他原因；但时间不允许，因为话题正催促进到祈祷的力量。正因如此，我确实提出了所有这些论点，好叫你们明白依撒格的祈祷如何解开了他妻子的不育；而且祈祷持续了那么久。经上说：\BibleQuote{依撒格不断为他的妻子黎贝加祈祷，天主垂允了他。}\BibleRef{创 25:21} 因为不要以为他一呼求天主就立刻蒙垂听；他曾在祈祷中花费了许多时间。如果你们想知道多少，我将确切地告诉你们这个数字。他在祈祷中花费了二十年。这从何显明？从上下文本身。因为圣经想要指出那个义人的信德、坚忍和智慧之爱，它没有中断甚至没有省略时间，而是也向我们表明，虽然隐晦，却是为了激发我们的懈怠；然而它并没有让时间不确定。那么请听它如何隐晦地向我们指明时间：\BibleQuote{依撒格娶黎贝加时，年四十岁。她是叙利亚人贝突耳的女儿。}\BibleRef{创 25:20} 你们听见他娶妻时多少岁：\BibleQuote{他娶黎贝加时年四十岁。}\BibleRef{创 25:20} 既然我们知道了娶妻时的年龄，让我们也学到他后来何时成为父亲，那时他多少岁，当他生了雅各伯时；我们就能看出他的妻子不育了多久；以及在那整个期间他继续向天主祈祷。那么他生雅各伯时多少岁？\BibleQuote{雅各伯出来时，手抓着哥哥的脚跟。因此给他起名叫雅各伯，那另一个叫厄撒乌。依撒格生他们时年六十岁。}\BibleRef{创 25:26} 所以如果娶黎贝加时他四十岁，生儿子时六十岁，显然他的妻子其间不育了二十年，而在这整个期间，依撒格继续向天主祈祷。\par

10. 此后，我们岂不感到羞愧、掩面，看到那个义人二十年坚持不懈、不放弃；而我们自己在一次或两次祈求后常常气馁和愤慨？然而，他确实对天主有极大的直言不讳，同样，他对恩赐的延迟没有感到不满，而是保持耐心，而我们则被无数罪孽重压，良心不安，对主没有善意；如果我们在说话之前没有被垂听，我们就迷惑，不耐烦地退缩，停止祈求——因此我们总是空手而归。谁曾为了一件事二十年恳求天主，如同这位义人一样？或者，谁只做了二十个月？昨天我还在说，有许多人祈祷时懈怠、打哈欠、伸懒腰、不断变换姿势，在祈祷中放纵一切疏忽——但今天我发现了另一种附着于他们祈祷的伤害，比那种更致命。因为许多人扑倒在地，以额头撞击地面，泪流满面，从心底痛苦呻吟，伸出双手，表现出极大的热忱，却用这种热情和激进反对自己的得救。因为他们恳求天主不是为自己的罪；也不是请求赦免他们所犯的过犯；而是完全将这种热忱用于反对他们的敌人，正像一个人磨快了剑，却不用武器对付敌人，而是刺穿自己的咽喉。所以这些人也利用祈祷不是为了赦免自己的罪，而是为了报复敌人；这无异于用剑刺向自己。这也是那恶者所设计的，好叫我们从各方面毁灭自己，无论是通过懈怠还是通过热忱。因为一类人因祈祷中的疏忽而激怒天主，通过懈怠表现出轻蔑；另一类人当他们表现出热忱时，却将热忱用于反对自己的得救。他说（魔鬼说）：\Quote{某个人懈怠：对我来说，这足以使他一无所获；这个人热忱且完全警醒；那么该做什么来达到同样结果？我不能削弱他的热忱，也不能使他陷入疏忽；我将以另一种方式图谋他的毁灭。如何？我将设法使他利用热忱去违犯法律。}（因为为自己的私仇祈祷是违犯法律。）\Quote{因此，他离去时不仅因热忱一无所获，而且还遭受了比懈怠所致的伤害更大的伤害。}这就是魔鬼的损害：他通过松懈毁灭一些人；通过热忱本身毁灭另一些人，当这热忱不是按照天主的法律而行时。\par

11. 但是，也值得听听他们祈祷的原话，以及这些话语是如何显出一种幼稚的心思，怎样一种婴孩的灵魂。我实在羞于重复它们，但又绝对必须重复它们，并且模仿那种粗俗的言辞。那么这些话是什么呢？\Quote{你为我报复我的仇敌，让他们也知道我有天主（站在我这边）。} 人哪，难道你们还不明白，当我们愤怒、恼怒和不耐烦时，我们并没有天主站在我们这边；而是当我们温柔、良善和谦和，并实践一切智慧之爱时，我们才有天主。同样，天主也曾说：\BibleQuote{你们的光要在人前照耀，好使他们看见你们的善行，光荣你们在天之父。}\BibleRef{玛 5:16} 你难道不明白，向天主祈求对付你的仇敌，是对天主的侮辱吗？怎么会是侮辱呢？有人会说。因为祂亲自说过：\BibleQuote{要为你们的仇敌祈祷；}并设立了这条神圣的律法。因此，当你要求立法者放宽他自己的律法，并呼召他制定与他自己相反的法律，并恳求那曾禁止你祈祷反对你的仇敌的主，听你祈祷反对你的仇敌；你这样做不是在祈祷，也不是呼召他，而是在侮辱立法者，并以醉酒的暴力对待他，而他本来必会赐予你祈祷所获得的好处。并且，请问，当你激怒那位必定会垂听的主时，又怎么可能在祈祷中被垂听呢？因为你这样做，是将自己的救恩推入深坑，冲下悬崖，就是当着君王的面击打你的仇敌。因为即使你不动手，你用言语击打他，这件事即使是对同为奴仆的，你也不敢做。至少在你长官面前你敢这样做吗？即使你曾立下无数功绩，你也必定立刻被带去处决。那么（我问你），在长官面前你不敢侮辱与你平等的人，但在天主面前这样做时，请问，你不战栗、不害怕吗？在恳求和祈祷的时候如此野蛮，把自己变成野兽，表现出比那个索要一百银元的人更无情吗？\BibleRef{玛 18:28} 因为你比他更傲慢，听听这故事本身：有一个人欠他的主人一万塔冷通，他没有钱偿还，就恳求主人忍耐，好让他变卖妻子、房屋和儿女，来清偿主人的债务。主人见他哀哭，就动了怜悯的心，免了他那一万塔冷通。他出去后，遇见一个欠他一百银元的同伴，便扼住他的喉咙，残忍无情地向他索讨。主人听说这事，就把他投入监狱，把先前免掉的那一万塔冷通的债务又加在他身上；他因对待同伴的残忍而受罚。\par

12. 现在你想想，你在某种意义上比那个人更麻木无情，你竟祈祷反对你的仇敌。他没有请求他的主人去索讨，而是自己索要那一百银元；而你竟呼求主人做这种无耻和禁止的索求。他是在外面，不在他主人眼前扼住同伴的喉咙；而你却在祈祷的当下，站在君王面前这样做。如果他这样做，既没有催促他的主人去索讨，而且事后出去，仍然没有得到宽恕；那么你，既煽动主人（去索取）这禁止的支付，又在他眼前这样做，你将要受什么惩罚呢？请告诉我。但是你的心因记恨而燃烧、膨胀，你的心升起，当在回忆中想起那伤害你的人时，你无法平息思想的膨胀。但是，要对抗这种炎症，要记起你自己所犯的罪，以及因将来的惩罚而产生的恐惧。回忆你向你的主人负责多少事，并要为所有那些事向他偿还；这种恐惧必然会克服那种愤怒，因为这实在比那激情更强大。在你祈祷的时候，要记起地狱、惩罚和报复；你将甚至不能把你的仇敌记在脑中。要使你的心怀痛悔，通过回忆你所犯的过犯而谦卑你的灵魂，愤怒甚至不能扰乱你。但这一切邪恶的原因在于，我们以极大的精确度审视所有人的罪，却以极大的松懈放过自己的罪。而我们应当相反——记住自己的过犯，永远不让别人的过犯进入思想。如果我们这样做，我们既会有天主慈爱，也会停止对邻人怀有永恒的愤怒，我们将永远不会有人作为仇敌；即使我们有时有，我们也会迅速结束他的敌意，并且获得对自己罪过的快速赦免。因为正如那珍藏对邻人伤害的记忆的人，不允许对自己罪的惩罚被消除；同样，那清除愤怒的人也会迅速清除自己的罪。因为如果我们这些邪恶、被情欲奴役的人，因天主的诫命而忽略别人对我们的一切过犯，那么那位爱人、良善、不受任何情欲影响的主，岂不更要忽略我们的过犯，以我们善待邻人的态度回报我们，赦免我们自己的罪吗？愿天主赐予我们达到此事，借着我们的主耶稣基督的恩宠和慈爱，愿光荣与权能归于祂，直到永远。阿们。\par

\SourceRecord{Translated by R. Blackburn.；Nicene and Post-Nicene Fathers, First Series；Vol. 9.；Edited by Philip Schaff.；Buffalo, NY: Christian Literature Publishing Co.,；1889.；<http://www.newadvent.org/fathers/1913.htm>.}
